
\section{Conclusion}

We studied why multimodal models favor text over vision when the two conflict, and found that the problem is not poor visual encoding. In many failure cases, models assign higher probability to the correct visual answer in intermediate layers, only to override it with the textual answer in deeper layers. We call this late-layer textual override. Importantly, the direction of prediction change distinguishes failures from successes: shifts toward text correlate with incorrect outputs, while shifts toward vision correlate with correct ones. This observation motivates CALRD, a training-free decoding method that restores intermediate predictions when override signatures are detected. Experiments on five MLLMs show 4--9\% improvements on conflict benchmarks without degrading performance on standard tasks. Our results suggest that improving reliability under conflict may require not better visual encoding, but better preservation of visual information through the forward pass.

\textbf{Limitations and future work.} Our work opens several avenues for future research. A natural next step is to examine how late-layer dynamics interact with model scale and pretraining data composition—questions that require controlled experiments with access to training corpora, beyond the scope of our mechanistic study. Additionally, exploring whether similar override patterns emerge across other modality pairs (\textit{e.g., }audio-text) may reveal general principles of multimodal integration. We view our mechanistic analysis as a foundation for such investigations.
